% Figure: residue clock mod 12, with +5 and +7 step arrows.
% Source for Figure~\ref{fig:clock}. TikZ only; regenerable by compiling main.tex.
\begin{tikzpicture}[
  nd/.style={circle,draw,minimum size=6mm,font=\footnotesize},
  stp/.style={->,thick,accent},
  lab/.style={font=\footnotesize\sffamily,text=slate},
]
\foreach \i in {0,...,11}
{
  \node[nd] (n\i) at ({90 - \i*30}:2.1) {\i};
}
% a +5 walk from 2: 2->7->0->5->10, then +7: 10->5
\draw[stp,bend left=18] (n2) to (n7);
\draw[stp,bend left=18] (n7) to (n0);
\draw[stp,bend left=18] (n0) to (n5);
\draw[stp,bend left=18] (n5) to (n10);
\draw[stp,dashed,bend right=18] (n10) to (n5);
\node[lab] at (0,-3.0) {$+5$ steps: $2\to7\to0\to5\to10$ (solid); then $+7$ (dashed): $10\to5$};
\node[lab] at (0,-3.4) {arithmetic on the circle is arithmetic on residues};
\end{tikzpicture}
